% University Project Report
\documentclass[12pt,a4paper]{article}
\usepackage[margin=1in]{geometry}
\usepackage{lmodern}
\usepackage[T1]{fontenc}
\usepackage{hyperref}
\usepackage{graphicx}
\usepackage{float}
\usepackage{booktabs}
\usepackage{listings}
\usepackage{xcolor}
\usepackage{tikz}
\usetikzlibrary{arrows.meta,positioning}

% Listings setup for C and Flex/Bison
\lstset{
  basicstyle=\ttfamily\footnotesize,
  keywordstyle=\color{blue}\bfseries,
  commentstyle=\color{gray}\itshape,
  stringstyle=\color{teal},
  numbers=left,
  numberstyle=\tiny\color{gray},
  frame=single,
  breaklines=true,
  tabsize=4,
  captionpos=b
}

\hypersetup{
  colorlinks=true,
  linkcolor=darkgray,
  urlcolor=blue,
  citecolor=blue
}

\title{\textbf{Compiler Design Project}\ \ \large Lexical, Syntax, and Code Generation Analysis}
\author{Roll: 2003001}
\date{}

\begin{document}
\maketitle
\tableofcontents
\newpage

\section{Introduction}
This report details the design and implementation of a compiler for a custom programming language. The language supports integer variables, basic arithmetic (addition and subtraction), assignments, and printing. The project transforms high-level source code into MASM-compatible assembly code through a four-stage pipeline: Lexical Analysis, Syntax Analysis, Intermediate Code Generation, and Target Code Generation.

The implementation uses standard compiler construction tools: \texttt{flex} for tokenization and \texttt{bison} for parsing, with backend logic implemented in C.

\paragraph{Problem Statement}
The compiler is designed to process the following snippet:
\begin{lstlisting}[language={}]
function main ()
start
       suppose A is an integer
       suppose B is an integer
       A equal 2
       B equal 3 add A sub 5
       print(B)
end
\end{lstlisting}

\section{System Architecture}
The compiler operates as a single-pass system where semantic actions in the parser trigger intermediate code generation.

\begin{figure}[H]
  \centering
  \resizebox{0.9\textwidth}{!}{%
    \begin{tikzpicture}[node distance=2.5cm, every node/.style={draw, rounded corners, minimum width=3cm, minimum height=1cm, align=center}, >={Stealth[length=3mm]}]
      \node (input) {Input Source\(.txt)};
      \node (lex) [right=of input] {Lexical Analysis\(Flex)};
      \node (syn) [right=of lex] {Syntax Analysis\(Bison)};
      \node (ir) [below=of syn] {Intermediate Code\(Struct Array)};
      \node (asm) [left=of ir] {Code Generation\(MASM Assembly)};
      \node (output) [left=of asm] {Output\(.asm / .exe)};

      \draw[->] (input) -- (lex);
      \draw[->] (lex) -- (syn);
      \draw[->] (syn) -- (ir);
      \draw[->] (ir) -- (asm);
      \draw[->] (asm) -- (output);
    \end{tikzpicture}%
  }
  \caption{Compiler Data Flow}
  \label{fig:pipeline}
\end{figure}

\section{Lexical Analysis}
The lexical analyzer scans the source code and groups characters into tokens. I used \texttt{flex} to define the rules.

\subsection{Token Specifications}
The lexer distinguishes between keywords (like \texttt{function}, \texttt{suppose}), operators, and literals. Note that in my implementation, the specific token names used in the C code (e.g., \texttt{ADDKW}, \texttt{STARTKW}) differ slightly from generic names to avoid conflicts.

\begin{table}[H]
  \centering
  \begin{tabular}{lll}
    \toprule
    \textbf{Pattern} & \textbf{Token Name} & \textbf{Semantic Value} \
    \midrule
    \texttt{function} & FUNCTION & - \
    \texttt{start} & STARTKW & - \
    \texttt{integer} & INTEGER & \texttt{yylval.int\_val = 1} \
    \texttt{add} & ADDKW & - \
    \texttt{[a-zA-Z][a-zA-Z0-9]*} & ID & \texttt{yylval.str\_val} \
    \texttt{[0-9]+} & ICONST & \texttt{yylval.int\_val} \
    \bottomrule
  \end{tabular}
  \caption{Key Token Definitions}
  \label{tab:tokens}
\end{table}

\subsection{Implementation (lexer.l)}
The \texttt{lexer.l} file defines these rules. I handle whitespace and newlines to track line numbers for error reporting.

... (FULL REPORT CONTINUES—REPLACED ALL "we" WITH "I")
